\section{Why ``Too Small'' $|B_{1,1}|$ Can Lead to Issues}
\label{sec:why-B1-small-issue}

\subsection{The Growth Rate Formula}

From Lemma~\ref{lem:r2-logratio-criterion-main}, the log-ratio growth rate is:
\[
\partial_t R_{12}(t,x_0) \ge (1-\rho)\,\xi_1(t)\,K_{\mu_0}(x_0,x_0)\,(-d_0(t))\,\frac{|B_{1,1}(t)|}{|f_1(t,x_0)|+\varepsilon} \ge 0
\]

The key observation is that the \textbf{lower bound on the growth rate is proportional to $|B_{1,1}(t)|$}.

\subsection{Why ``Too Small'' $|B_{1,1}|$ Is Problematic}

\paragraph*{1. Insufficient Growth Rate}

If $|B_{1,1}(t)|$ is very small (e.g., $|B_{1,1}(t)| \ll \varepsilon$), then:
\[
\partial_t R_{12}(t,x_0) \ge (1-\rho)\,\xi_1(t)\,K_{\mu_0}(x_0,x_0)\,(-d_0(t))\,\frac{|B_{1,1}(t)|}{|f_1(t,x_0)|+\varepsilon} \approx 0
\]

The growth rate becomes \textbf{extremely small}, meaning:
\begin{itemize}
\item The log-ratio $R_{12}$ grows very slowly (or barely at all)
\item Dominance is established very slowly, if at all
\item The channel might not effectively learn the target function
\end{itemize}

\paragraph*{2. Numerical Errors and Discretization}

In practice, we use discrete-time updates (gradient descent with step size $\epsilon$). If the continuous-time growth rate is too small:
\begin{itemize}
\item \textbf{Discretization error} $O(\epsilon)$ might overwhelm the growth $O(|B_{1,1}| \cdot \epsilon)$
\item If $|B_{1,1}| \ll 1/\epsilon$, then discretization error dominates
\item The actual discrete-time dynamics might not preserve dominance
\item Numerical round-off errors can cause the sign of $B_{1,1}$ to flip, violating the assumption
\end{itemize}

\paragraph*{3. Sign Coherence Violation}

The lemma requires that \textbf{$B_{1,1}(t)$ has the same sign as $f_1(t,x_0)$}. If $|B_{1,1}|$ is too small:
\begin{itemize}
\item Small perturbations (numerical errors, noise) can flip the sign
\item The assumption ``$B_{1,1}$ has the same sign as $f_1$'' might be violated
\item The lower bound $\partial_t R_{12} \ge 0$ might not hold
\item Dominance could actually \textbf{decrease} instead of increase
\end{itemize}

\paragraph*{4. Ineffective Learning}

$B_{1,1}(t)$ is the backpropagated signal that drives learning for channel 1. If it's too small:
\begin{itemize}
\item The channel receives very weak gradient signals
\item The update $\partial_t f_1 \propto -B_{1,1}$ is very small
\item The channel learns very slowly or not at all
\item The dominance might not be ``meaningful'' (the channel isn't actually learning)
\end{itemize}

\paragraph*{5. Competition from Other Channels}

Even if the growth rate is positive, if it's too small:
\begin{itemize}
\item Other channels (with larger $|B_k|$) might catch up or overtake
\item The stability assumption \eqref{eq:r2-stability-assumption-main} might be violated if $B_2$ grows faster
\item Channel switching could occur, breaking the specialization
\end{itemize}

\subsection{Quantitative Threshold}

A reasonable condition for ``not too small'' would be:
\[
|B_{1,1}(t)| \ge C \cdot \max\left(\varepsilon, \frac{\epsilon}{\xi_1(t) K_{\mu_0}(x_0,x_0) (-d_0(t))}\right)
\]
for some constant $C > 0$, where:
\begin{itemize}
\item $\varepsilon$ is the regularization parameter in the log-ratio
\item $\epsilon$ is the discretization step size (if using discrete-time)
\item The second term ensures discretization error doesn't dominate
\end{itemize}

\subsection{Connection to the Stability Assumption}

The stability assumption \eqref{eq:r2-stability-assumption-main} requires:
\[
|B_{2,1}(t)| \le \rho \frac{|f_2(t,x_0)|+\varepsilon}{|f_1(t,x_0)|+\varepsilon} |B_{1,1}(t)|
\]

If $|B_{1,1}|$ is too small, then:
\begin{itemize}
\item $|B_{2,1}|$ is also very small (bounded by $\rho |B_{1,1}|$)
\item Both channels receive weak signals
\item Neither channel learns effectively
\item Specialization might not occur
\end{itemize}

\subsection{Summary}

The phrase ``whenever $|B_{1,1}|$ is not too small'' means:
\begin{enumerate}
\item \textbf{Growth rate is meaningful:} The log-ratio grows at a rate that's not overwhelmed by numerical errors
\item \textbf{Sign coherence is stable:} Small perturbations don't flip the sign of $B_{1,1}$ relative to $f_1$
\item \textbf{Effective learning:} The channel receives strong enough gradient signals to actually learn
\item \textbf{Dominance is established:} The growth is fast enough to establish and maintain dominance before other effects interfere
\end{enumerate}

Without this condition, the lemma still guarantees $\partial_t R_{12} \ge 0$ (non-decreasing), but the growth might be so slow that it's practically ineffective, or numerical errors might cause violations of the assumptions.
